%\documentclass[twocolumn]{article}\
\documentclass{article}
\usepackage{amsmath, amssymb, cancel, mathtools, bm}
\usepackage[left=.5in, right=.5in, top=1in, bottom=1in]{geometry}
\usepackage[most]{tcolorbox}
\usepackage[skip=1em,indent=0pt]{parskip}
\setlength{\parindent}{0pt}
\newcommand{\svec}[1]{\underset{\sim}{\bm{#1}}} % Vector symbol (tilde under vector)
\DeclareMathOperator{\EX}{\mathbb{E}} % Expected Value symbol
\DeclareMathOperator{\ext}{\EX_{\theta}} % Expected Value symbol
\DeclareMathOperator{\vart}{Var_{\theta}} % Expected Value symbol
\newcommand{\indep}{\perp\!\!\!\!\perp} % Independence symbol
\newcommand{\real}{\mathbb{R}} % Simplified 'Reals' indicator
\newcommand{\pto}{\overset{P}{\to}}
\newcommand{\asto}{\overset{a.s.}{\to}}
\newcommand{\rthto}{\overset{L^r}{\to}}
\newcommand{\dto}{\overset{D}{\to}}
\newcommand{\xtb}{x^{\top}\svec{\beta}}
\newcommand{\xitb}{x_i^{\top}\svec{\beta}}
% Force all aggregate symbols to always put values above/below
\let\Oldint=\int
\let\Oldsum=\sum
\let\Oldprod=\prod
\let\Oldbigcup=\bigcup
\let\Oldbigcap=\bigcap
\let\Oldlim=\lim
\renewcommand{\int}{\Oldint\limits} 
\renewcommand{\sum}{\Oldsum\limits} 
\renewcommand{\prod}{\Oldprod\limits} 
\renewcommand{\bigcup}{\Oldbigcup\limits} 
\renewcommand{\bigcap}{\Oldbigcap\limits} 
\renewcommand{\lim}{\Oldlim\limits} 
\newcommand{\infint}{\int_{-\infty}^{\infty}}
\newcommand{\iiddist}{\overset{\mathrm{iid}}{\sim}}
\newcommand{\inddist}{\overset{\mathrm{ind}}{\sim}}
\newcommand{\norm}[1]{\left\lVert#1\right\rVert}
\newcommand{\boldred}[1]{\textbf{\textcolor{red}{#1}}}
\usepackage[linktoc=all]{hyperref}



\begin{document}

\tableofcontents

\newpage

%%%%%%%%%%%%%%%%%%%%%%%%%%%%%%%%%%%%%%%%%%%%%%%%%%%%%%%%%%%%%%%%%%%%%%%%%%%
\section{GLM Overview}

\subsection{Example}
Normal Regression: $y_i = \beta_0 + \beta_1 x_i + \epsilon_i$: $\epsilon_i \iiddist N(0, \sigma^2)$

Logistic Regression: $y_i \overset{ind}{\sim} Bern(p_i)$: $logit(p_i) = \beta_0 + \beta_1 x_i$, $ log(\frac{p_i}{1-p_i})$

Differences: 
\begin{itemize}
    \item Support of $y_i$: $y_i \in \real$ vs. $y_i \in \{0, 1\}$
    \item Link functions: $g(\mu_i) = \mu_i$ vs. $g(p_i) = log(\frac{p_i}{1-p_i})$
    \item Likelihoods: $y_i \overset{ind}{\sim} N(\mu_i, \sigma^2)$ vs. $y_i \overset{ind}{\sim} Bern(p_i)$
\end{itemize}

Similarities:
\begin{itemize}
    \item Univariate models
    \item Exponential families
    \item Indepdently distributed
    \item Linear (or otherwise systematic) relationship with a function of $E(y_i)$
    \item Have link functions
    \item Assuming that $y_i$ is random
    \item $\beta$'s are unknown
\end{itemize}
    

\subsection{GLMs}
Generalized Linear Model:

\begin{itemize}
    \item Random component: $Y_i \sim f(\mu_i)$
    \item Link component: $g(\mu_i) = \eta_i$
    \item Systematic component: $\eta_i = \xitb$
\end{itemize}

We generally use LMs for the following:
\begin{itemize}
    \item Inference: understand relationships
    \item Prediction: hevaior of specific x's
\end{itemize}


\subsection{Poisson Regression}

$Y_i \inddist Poisson(\mu_i), log(\mu_i) = x_i^{\top}\beta$

Assumptions: 
\begin{itemize}
    \item Linearity
    \item Independence
    \item Poisson distribution on response var
\end{itemize}


\subsection{Exponential Family Calculations}

With the Exponential Family log-likelihood written as: $l(\theta; y) = w(\theta)t(y) + ln(a(\theta)) + ln(b(y))$ where $ln(a(\theta)) = c(\theta)$ and $ln(b(y)) = d(y)$:

$Score = w'(\theta)t(y) + c'(\theta)$

$\EX[t(y)] = -\frac{c'(\theta)}{w'(\theta)}$

$Var(t(y)) = \frac{w''(\theta)c'(\theta) - w'(\theta)c''(\theta)}{(w'(\theta)^3)}$

\underline{Example}: Poisson: $f(y; \lambda) = \frac{\lambda^y e^{-\lambda}}{y!} = exp\{yln(\lambda - \lambda - ln(y!))\}$

In this case, $c(\lambda) = -\lambda$ and $w(\lambda) = ln(\lambda)$

Sufficient Statistic: $t(y) = y$

$w'(\lambda) = \frac{1}{\lambda}, w''(\lambda) = -\frac{1}{\lambda^2}, c'(\lambda) = -1, c''(\lambda) = 0$

Mean: $\EX[t(y)] = \EX[y] = -\frac{c'(\lambda)}{w'(\lambda)} = -\frac{1}{\frac{-1}{\lambda}} = \lambda$

Variance: $Var(t(y)) = Var(y) = \frac{w''(\lambda)c'(\lambda) - w'(\lambda)c''(\lambda)}{(w'(\lambda))^3} = \frac{(-1)(-\frac{1}{\lambda^2}) - 0}{\frac{1}{\lambda^3}} = \frac{\frac{1}{\lambda^2}}{\frac{1}{\lambda^3}} = \lambda$

Score: $w'(\theta)t(y) + c'(\theta) = (\frac{1}{\lambda})y + (-1)$. Verify $\EX[Score] = 0: \EX[\frac{1}{\lambda}y-1] = \frac{1}{\lambda}\EX[y] - 1 = \frac{\lambda}{\lambda} - 1 = 0$

\subsection{Link Function}

$\eta_i = g(\mu_i)$

Requirements (all with respect to the mean $\mu$, not $x$):
\begin{itemize}
    \item One-to-One ($g(a) = g(b) \implies a=b$)
    \item Monotone (Monotone increasing means $a >b \implies g(a) > g(b)$)
    \item Continuous and differentiable ($lim_{\mu \to c}g(\mu) = g(c)$)
\end{itemize}

Canonical Link: $g*(w(\theta_i)) = w(\theta_i)$

Example, 

Binomial: $w(\pi_i) = log(\frac{\pi_i}{1-\pi_i}) \implies g(\pi_i) = log(\frac{\pi_i}{1-\pi_i}), g*(w(\pi_i)) = w(\pi_i)$


\section{}

With $U$ as the score function, and $\mathbb{J}$ as the variance:

$U \sim N(\svec{0}, \mathbb{J})$

$\mathbb{J}^{-1/2} U \sim N(0,1)$

$(\mathbb{J}^{-1/2} U)^{\top}\mathbb{J}^{-1/2} U = U^{\top}\mathbb{J}^{-1}U \sim \chi^2(k)$


\section{Prediction Intervals}

Using regular process for confint on predictions of using the variance of $\frac{X \mathbb{I}X^{\top}}{g^{-1}(\hat{\mu})}$ provides too large of an interval. The better method is to get interval on $\hat{\eta}$ then transform bounds to $\hat{\mu}$:

$\hat{\eta} = x\hat{\beta} \sim N(x\beta, X Cov(\hat{\beta})X^{\top})$ Use that variance to get $SE(\hat{\eta})$

$\hat{\eta} \pm 1.96 SE(\hat{\eta}) = [\hat{\eta}_1, \hat{\eta}_2], \hat{\mu}_2 = g^{-1}(\hat{\eta}_2)$

\section{Residuals}

Perason Residuals: $r_{p_i} = \frac{y_i - \hat{\mu}_i}{\sqrt{Var(Y_i)}}$

Examples: 

Poisson: $\frac{y_i - \hat{\mu}_i}{\sqrt{\hat{\mu}_i}}$

Logit: $\frac{y_i - \hat{\mu}_i}{}$

Anscombe: R doesn't really even touch

Deviance: uses the individual observation log-likelihood

Working: Similar to the Z statistic (difference from mean divided by SE)

\section{Picking Link Functions}

There are lots of ways to pick link functions:
\begin{itemize}
    \item Interpretation
    \item AIC/BIC
    \item Specificity
    \item etc.
\end{itemize}

In some cases, the link function doesn't make much of a difference. However, sometimes it does. 

\section{Crabs Analysis}

Model with 'is.very.dark' as the covariate. 2 common ways of writing:

1. 

$\eta_{ij} = \alpha_0 + \gamma_j$ : $Y_{ij} = $ whether or not the ith crab in the jth color group has a sateliitle ($j  = 1, 2; i = 1, ..., n$)

Often constrained such that $\eta_{i1} = \alpha_0 + \gamma_1, \eta_{i2} = \alpha_0 + \gamma_2, \gamma_1 + \gamma_2 = 0$ ("sum to zero" constraint)

ie. $x = \begin{pmatrix}
    1 & -1 \\ 1 & -1 \\ 1 & 1 \\ 1 & 1 \\ \vdots
\end{pmatrix}$

2. 

$\eta_i = \beta_0 + I(x_i = True)\beta_1$

When written with the indicator, that's the same as constraint $\gamma_1 = 0$

ie. $x = \begin{pmatrix}
    1 & 0 \\ 1 & 0 \\ 1 & 1 \\ 1 & 1 \\ \vdots
\end{pmatrix}$


Note:

$\hat{\beta}_0 = \hat{\alpha}_0 + \hat{\gamma}_1$

$\hat{\beta}_0 + \hat{\beta}_1 = \hat{\alpha}_0 + \hat{\gamma}_2$

$\beta_0$ is the log odds that a not very dark crab has a satellite

$e^{\beta_0}$ is the odds that a not very dark crab has a satellite

$\beta_0 + \beta_1$ is the log odds that a very dark crab has a satellite

$e^{\beta_0 + \beta_1}$ is the odds that a very dark crab has a satellite


$\pi_{i1}$ is the probability of satellite for not very dark

$\pi_{i2}$ is the probability of satellite for very dark

$ln(\frac{\pi_{i2}}{(1-\pi_{i2})}) = \beta_0 + \beta_1$

$ln(\frac{\pi_{i1}}{(1-\pi_{i1})}) = \beta_0$

$ln(\frac{\pi_{i2}/(1-\pi_{i2})}{\pi_{i1}/(1-\pi_{i1})}) = \beta_1$ which is the log odds ratio that a crab has a satellite for a very dark vs. not very dark

Now with a 4-level factor as the covariate

$\eta_i = \beta_0 + I(x_i="DarkMed")\beta_{dm} + I(x_i="LightMed)\beta_{lm} + I(x_i="Med")\beta_{med}$

$H_0: \beta_{dm} = \beta_{lm} = \beta_{med} = 0, H_1:$ at leat one is not 0

Test with Delta Deviance and compare to $\chi^2(3)$. We get a p-value of .0033 so we rejec the $H_0$ and conclude there is enough evidence that at least of the $\beta$s should be in the model. 

For testing a specific $\beta$:

$H_0: \beta_{dm} = 0, H_1: \beta_{dm} \neq 0$. p-value is 0.04. Any significance less than this would find that dm is not different from dark

Now, for testing contrasts between $\beta$s:

$H_0: \beta_{med} - \beta_{lm} = 0, H_1: \beta_{med} - \beta_{lm} \neq 0$

Contrast: $c = (0, 0, -1, 1)^{\top}$, usual contrast rules apply

$e^{\beta_{med}}$ is the odds ratio that a crab has a satellite for mdeium vs. dark crabs ($\eta_{i,med} - \eta_{i,d} = \beta_{med}, log(\frac{\pi_{1,med}/(1-\pi_{i,med})}{\pi_{1,d}/(1-\pi_{i,d})})$)

\section{4.5} 

Final model from selection:

$y_i \inddist Bern(\pi_i)$

$ln(\frac{\pi_i}{1-\pi_i}) = \beta_0 + \beta_1x_{i1} + \beta_2x_{i2}$

$x_{i1}$ represents the width of crab i, $x_{12}$ represents the indicator for 'is very dark' for crab i

For a 1 unit increase in width, the log-odds increase by $\beta_1$.

The odds change by a factor of $e^{\beta_1}$ when going from one width to a width of 1 unit more.

$e^{\beta_1}$ is the odds-ratio of having a satellite for two crabs one unit width difference. "Reversing" it, $e^{\beta_1}$ is the odds-ratio of a crab being one unit wider for a crab that does or does not have a satellie.

\end{document}